\documentclass[10pt,twoside]{article}
\newcommand{\BedrockArchitectureName}{Bedrock}
\newcommand{\BedrockDocumentTitle}{Bedrock ELF ABI}
\newcommand{\BedrockRunningTitle}{ELF ABI}
\newcommand{\BedrockFooterTitle}{BEDROCK ARCHITECTURE}
\usepackage{bedrock-reference}
\begin{document}
\BedrockMakeTitlePage{Bedrock ELF ABI}{ELF Object, Linking, and Loading ABI}{Draft}
\BedrockDocumentHasFiguresfalse
\BedrockFrontMatter

\clearpage
\section{Scope}\label{page:scope}
This document defines the Bedrock binary-format contract for ELF objects,
linking, relocation, loading, dynamic linking, TLS metadata, and code models.
The Bedrock ISA Specification supplies the referenced architectural objects.
Language bindings and application interfaces each own their higher-level
contracts.

\subsection{Contract Position}
\BedrockContractDiagram{binary}

\subsection{Profile}
\BedrockField{Profile:}{\texttt{bedrock-elf}}

\clearpage
\section{Terminology}\label{page:terminology}
These terms are used by the Bedrock ELF ABI. ISA-level address, register, and effective-address terms keep the meanings defined by the architecture manual.
\begin{description}[style=nextline,leftmargin=1.45in,labelwidth=1.35in,itemsep=3pt,topsep=3pt]
\item[ELF ABI] The ELF, relocation, dynamic-linking, TLS, and code-model contract shared by language bindings and loaders.
\item[pre-segment virtual address] The stable coordinate used by ELF symbols, relocations, ordinary function addresses, and ordinary data pointers before segment pre-translation.
\item[load base] The runtime address bias applied to a loaded ET\_DYN object before its symbol values and relative relocations are interpreted.
\item[ABI-visible linkage domain] The region in which ordinary ELF code addresses are near-call entry points represented solely by pre-segment coordinates.
\item[near-call entry address] A pre-segment code address that can be reached with ordinary CALL and returned from with ordinary RET inside the current linkage domain.
\item[veneer] A near-call entry point supplied by a linker or loader to encapsulate a more complex transfer sequence for ordinary ABI callers.
\item[code model] A set of placement and relocation assumptions used by code generation to choose direct displacement forms, GOT/PLT forms, or full-width address materialization.
\end{description}

\clearpage
\section{ELF Object Format}\label{page:elf-object-format}
Bedrock objects use the 64-bit little-endian ELF format and machine identifier
\texttt{EM\_\allowbreak{}BEDROCK} (\texttt{0xffb0}). A producer may
emit relocatable objects, executable images, or shared/position-independent
images using \texttt{ET\_REL}, \texttt{ET\_EXEC}, or \texttt{ET\_DYN}.

All relocation sections use \texttt{Elf64\_Rela} entries with explicit
addends and the \texttt{.rela} section-name stem. Ordinary symbol and string
tables use \texttt{.symtab} and \texttt{.strtab}. The baseline value of
\texttt{e\_flags} is zero; a consumer accepts zero as the baseline profile and
rejects every unknown nonzero value.

\subsection{ELF Identification Bytes}
\BedrockTableCaption{ELF e\_ident Values}
\begingroup\footnotesize
\setlength{\tabcolsep}{2pt}
\begin{longtable}{@{}p{1.55in}p{3.95in}@{}}
\toprule
\textbf{e\_ident Field} & \textbf{Value}\\
\midrule
\endhead
\texttt{EI\_\allowbreak{}MAG} & \texttt{0x7f 'E' 'L' 'F'}\\
\texttt{EI\_\allowbreak{}CLASS} & \texttt{ELFCLASS64}\\
\texttt{EI\_\allowbreak{}DATA} & \texttt{ELFDATA2LSB}\\
\texttt{EI\_\allowbreak{}VERSION} & \texttt{EV\_\allowbreak{}CURRENT}\\
\texttt{EI\_\allowbreak{}OSABI} & \texttt{ELFOSABI\_\allowbreak{}NONE}\\
\texttt{EI\_\allowbreak{}ABIVERSION} & \texttt{0}\\
\bottomrule
\end{longtable}
\endgroup
\subsection{ELF Header Sizes}
\BedrockTableCaption{ELF Header Size Values}
\begingroup\footnotesize
\setlength{\tabcolsep}{2pt}
\begin{longtable}{@{}p{1.55in}p{3.95in}@{}}
\toprule
\textbf{Header Field} & \textbf{Bytes}\\
\midrule
\endhead
\texttt{e\_\allowbreak{}ehsize} & 64\\
\texttt{e\_\allowbreak{}phentsize} & 56\\
\texttt{e\_\allowbreak{}shentsize} & 64\\
\bottomrule
\end{longtable}
\endgroup
\subsection{Entry Point Semantics}
\BedrockTableCaption{ELF Entry Point Rules}
\begingroup\footnotesize
\setlength{\tabcolsep}{2pt}
\begin{longtable}{@{}p{1.75in}p{3.75in}@{}}
\toprule
\textbf{File Type} & \textbf{e\_entry Meaning}\\
\midrule
\endhead
\texttt{ET\_\allowbreak{}REL} & 0\\
\texttt{ET\_\allowbreak{}EXEC} & absolute pre-segment virtual byte address\\
\texttt{ET\_\allowbreak{}DYN (executable)} & load-bias-relative pre-segment virtual byte address; runtime entry is load base + e entry\\
\texttt{ET\_\allowbreak{}DYN (shared object, entry value 0)} & 0\\
\bottomrule
\end{longtable}
\endgroup
\subsection{Program Header Policy}
The standard \texttt{PT\_LOAD}, \texttt{PT\_DYNAMIC},
\texttt{PT\_INTERP}, \texttt{PT\_NOTE}, and \texttt{PT\_TLS} program
headers retain their ordinary ELF meanings.

Program permission bits use the standard values \texttt{PF\_X=1},
\texttt{PF\_W=2}, and \texttt{PF\_R=4}. Every \texttt{PT\_LOAD} has
\texttt{p\_align} equal to the selected page size or a larger power-of-two
alignment; the minimum supported page alignment is 4096 bytes.

\subsection{Section and Symbol Policy}\label{section:section-and-symbol-policy}
Standard ELF section types and the \texttt{SHF\_WRITE},
\texttt{SHF\_ALLOC}, \texttt{SHF\_EXECINSTR}, and \texttt{SHF\_TLS} flags
retain their generic meanings. Symbol binding and type occupy
\texttt{st\_info} in the ordinary ELF encoding. Architecture-specific
\texttt{st\_other} bits are reserved and must be zero.

In an \texttt{ET\_REL} object, \texttt{st\_value} is a byte offset from the
start of the defining section. After final link, it is the symbol's
pre-segment virtual byte address. Ordinary function symbols retain their
pre-segment address representation, and ordinary object symbols retain their
byte-address representation.

\clearpage
\section{Program Loading}\label{page:program-loading}
Bedrock ELF program loading follows the ordinary ELF segment model. This ELF
ABI defines executable mapping, loader-visible object rules, and entry-image
state.
\subsection{Loading Rules}
The loader may execute \texttt{ET\_EXEC} and executable \texttt{ET\_DYN}
images. It honors \texttt{PT\_INTERP}, maps each loadable segment at its
declared page alignment, and rejects an alignment below 4096 bytes. For a
segment whose \texttt{p\_memsz} exceeds \texttt{p\_filesz}, the loader fills
every byte in the half-open range
$[\texttt{p\_vaddr}+\texttt{p\_filesz},
\texttt{p\_vaddr}+\texttt{p\_memsz})$ with zero. Object and section alignment
requirements remain binding. The zero-filled gap may begin at any byte
alignment.
\subsection{Page Permissions}
Every \texttt{PT\_LOAD} shall set \texttt{PF\_R}. A loader rejects a
\texttt{PT\_LOAD} whose flags are zero or pair \texttt{PF\_W} or
\texttt{PF\_X} with a clear \texttt{PF\_R}. For an accepted segment, its flags
select one Normal leaf AM as shown below. The loader rejects
\texttt{PF\_R|PF\_W|PF\_X} because the available leaf AMs separate writable
and executable mappings.

\BedrockTableCaption{ELF Segment Permission Mapping}
\begingroup\footnotesize
\setlength{\tabcolsep}{2pt}
\begin{longtable}{@{}p{1.2in}p{1.45in}p{2.85in}@{}}
\toprule
\textbf{ELF Flags} & \textbf{Leaf AM and Fields} & \textbf{Meaning}\\
\midrule
\endhead
\texttt{PF\_\allowbreak{}R} & \texttt{AM=R}, P, U & readable user page\\
\texttt{PF\_\allowbreak{}R\textbar{}PF\_\allowbreak{}W} & \texttt{AM=RW}, P, U & readable writable user page\\
\texttt{PF\_\allowbreak{}R\textbar{}PF\_\allowbreak{}X} & \texttt{AM=RX}, P, U & readable executable user page\\
\bottomrule
\end{longtable}
\endgroup
\subsection{Execution Environment}
The ELF contract interprets \texttt{e\_entry}, \texttt{p\_vaddr}, symbol
values, relocation places, and linked code and data addresses in one stable
pre-segment virtual-address coordinate. Every ABI-permitted path that consumes
one such coordinate preserves its arithmetic and linkage meaning and resolves
an actual access to the intended byte of the intended object. At transfer to \texttt{e\_entry},
the code context must permit fetch from the entry address, the data context
must cover the mapped ordinary data regions, and the stack context must permit
reads and writes to the entry stack.

The loader maps each post-segment linear range required by these coordinates
to the intended physical pages with the AM and U permissions required by the
ELF image. Process entry, dynamic linking, module transitions, TLS paths,
unwind and debugger state, JIT publication, and shared pointer conventions
preserve the same coordinate invariants.

Executable images and ordinary shared objects form one ABI-visible near-call
linkage domain. Their function symbols use pre-segment code addresses. The
ordinary dynamic-linking contract covers near calls, and a linker-supplied
veneer brings a segment-changing transfer into that domain.

\subsection{Program Entry State}
Before transferring control, the loader completes image relocations and
publishes executable bytes. At transfer, \texttt{PC} equals the run-time value
of \texttt{e\_entry}; \texttt{SP} names a writable entry stack with 16-byte
alignment; and the code, data, and stack segment contexts provide the image
described above. A TLS image uses \texttt{GS0} as its thread-pointer base.

The loader clears \texttt{FLAGS}, \texttt{FSTATUS}, \texttt{FFLAGS}, and
\texttt{VSTATUS}.
General, floating-point, vector, and predicate registers outside the stated
entry contract carry unspecified values. An external process-entry ABI defines
the entry-stack payload and the interpretation of initial arguments. This
language-independent state permits an ELF entry point to establish the
run-time convention selected by its implementation language.

(:elf-entry-state:)

\clearpage
\section{Sections and Symbols}\label{page:sections-and-symbols}
Bedrock uses ordinary ELF sections with the usual ELF meanings.
\subsection{Symbol Model}
Bedrock supports the standard local, global, and weak bindings and the
default, hidden, and protected visibility classes. The standard notype,
object, function, TLS, ifunc, section, and file symbol types retain their ELF
meanings.

\begin{itemize}
\item Function symbols name the first instruction of the callable entry point.
\item Object symbols name the first byte of the object storage.
\item Symbol values follow the Section and Symbol Policy in
Section~\ref{section:section-and-symbol-policy}.
\end{itemize}
\clearpage
\section{Relocation Set}\label{page:relocation-set}
Bedrock ELF objects use RELA relocation entries with explicit addends. The Bedrock ABI defines how each relocation type patches instruction or data fields.
\subsection{Expression Model}
Every static relocation starts with the canonical expression
\texttt{symbol + addend},
where the symbol and explicit RELA addend are interpreted in the address space
of the relocation family. Ordinary relocations use ABI symbol addresses, and
TLS relocations use GS0-relative offsets.

Absolute and section-relative results are permitted when the selected
relocation type defines them. PC-relative address materialization subtracts
\texttt{place}, the address of the patched field. Direct branch and call payloads
subtract \texttt{next\_pc}, the address of the instruction following the patched
instruction. An out-of-range result for the signedness and width stated in the
relocation table is a link error, and the linker rejects the relocation.

For a PC-relative address-materialization relocation, let \texttt{I} be the
address of byte zero of the containing instruction and let
$\delta=\texttt{place}-\texttt{I}$ be the relocation field's byte offset. Let
$A_{\mathrm{src}}$ be the addend in the source expression before field-offset
bias. The encoded value shall equal the target plus $A_{\mathrm{src}}$ minus
\texttt{I}. Because the relocation-table expressions subtract \texttt{place}, the
assembler records the ELF addend as
$\texttt{addend}=A_{\mathrm{src}}+\delta$ for
\texttt{PCREL}, \texttt{GOTPCREL}, \texttt{GOT\_BASE\_PCREL}, and
\texttt{TLSDESC\_GOTPCREL} relocations. The assembler supplies every field
offset through the addend. The PLT addend described in Section~7.3 is one instance of this
general rule.
\subsection{Relocation Metasyntax Vocabulary}
\BedrockTableCaption{Relocation Metasyntax Vocabulary}
\begingroup\footnotesize
\setlength{\tabcolsep}{2pt}
\begin{longtable}{@{}p{1.45in}p{4.05in}@{}}
\toprule
\textbf{Identifier or function} & \textbf{Meaning}\\
\midrule
\endhead
\texttt{symbol} & Runtime byte address of the symbol named by the relocation after symbol resolution.\\
\texttt{addend} & Explicit addend stored in the ELF64 Rela entry.\\
\texttt{place} & Runtime byte address of the relocation field being patched.\\
\texttt{next\_pc} & Runtime byte address of the next instruction after the instruction being patched.\\
\texttt{section\_base} & Runtime byte address of the beginning of the symbol's defining section.\\
\texttt{symbol\_size} & Byte size of the resolved symbol.\\
\texttt{load\_base} & Runtime load base address of the loaded object that owns the relocation.\\
\texttt{got\_base} & Runtime byte address of the global offset table base for the loaded object.\\
\texttt{got(symbol)} & Runtime byte address of the GOT slot allocated for \texttt{symbol}.\\
\texttt{plt(symbol)} & Runtime byte address of the PLT entry allocated for \texttt{symbol}.\\
\texttt{tls(symbol)} & Byte offset of the TLS symbol from the ABI TLS base in GS0.\\
\texttt{tlsdesc(symbol)} & Runtime byte address of the TLSDESC entry allocated for the TLS symbol.\\
\texttt{resolver(x)} & Runtime byte address returned by calling the resolver function at address \texttt{x}.\\
\texttt{copy(symbol, symbol\_size)} & Copy operation over the resolved symbol's bytes.\\
\texttt{tls\_descriptor(symbol)} & Adjacent function and argument words selected for the TLS symbol.\\
\bottomrule
\end{longtable}
\endgroup
\subsection{GOT and PLT Model}
The ordinary GOT stores 64-bit addresses for position-independent code and
dynamic binding. A GOT relocation allocates or references the symbol's slot in
the loaded object. The PLT contains executable near-call entries; a PLT
relocation names that entry as its callable target. Every
ordinary PLT slot is resolved before application entry or constructors execute.

GOT and PLT function addresses stay inside the ABI-visible near-call linkage
domain.
\subsection{Relocations}
\BedrockTableCaption{Bedrock ELF Relocation Types}
\begingroup\footnotesize
\setlength{\tabcolsep}{2pt}
\begin{longtable}{@{}p{0.45in}p{2.15in}p{0.85in}p{2.05in}@{}}
\toprule
\textbf{Type} & \textbf{Name} & \textbf{Size} & \textbf{Calculation}\\
\midrule
\endhead
(:elf-relocations:base.relocations.R_BEDROCK_NONE,base.relocations.R_BEDROCK_ABS8,base.relocations.R_BEDROCK_ABS16,base.relocations.R_BEDROCK_ABS32S,base.relocations.R_BEDROCK_ABS64,base.relocations.R_BEDROCK_IMM8S,base.relocations.R_BEDROCK_IMM16S,base.relocations.R_BEDROCK_IMM32S,base.relocations.R_BEDROCK_IMM64,base.relocations.R_BEDROCK_DISP8S,base.relocations.R_BEDROCK_DISP16S,base.relocations.R_BEDROCK_DISP32S,base.relocations.R_BEDROCK_DISP64,base.relocations.R_BEDROCK_PCREL8S,base.relocations.R_BEDROCK_PCREL16S,base.relocations.R_BEDROCK_PCREL32S,base.relocations.R_BEDROCK_PCREL64,base.relocations.R_BEDROCK_BRDISP8S,base.relocations.R_BEDROCK_BRDISP16S,base.relocations.R_BEDROCK_BRDISP32S,base.relocations.R_BEDROCK_CALL16S,base.relocations.R_BEDROCK_CALL32S,base.relocations.R_BEDROCK_SECTION_REL32,base.relocations.R_BEDROCK_SECTION_REL64,base.relocations.R_BEDROCK_GOT64,base.relocations.R_BEDROCK_GOTPCREL32S,base.relocations.R_BEDROCK_GOTPCREL64,base.relocations.R_BEDROCK_GOTOFF32S,base.relocations.R_BEDROCK_GOTOFF64,base.relocations.R_BEDROCK_GOT_BASE_PCREL32S,base.relocations.R_BEDROCK_GOT_BASE_PCREL64,base.relocations.R_BEDROCK_PLT16S,base.relocations.R_BEDROCK_PLT32S,base.relocations.R_BEDROCK_PLT64,base.relocations.R_BEDROCK_RELATIVE,base.relocations.R_BEDROCK_GLOB_DAT,base.relocations.R_BEDROCK_JUMP_SLOT,base.relocations.R_BEDROCK_COPY,base.relocations.R_BEDROCK_IRELATIVE,base.relocations.R_BEDROCK_TLS_OFFSET32S,base.relocations.R_BEDROCK_TLS_OFFSET64,base.relocations.R_BEDROCK_TLSDESC_GOTPCREL32S,base.relocations.R_BEDROCK_TLSDESC_GOTPCREL64,base.relocations.R_BEDROCK_TLSDESC_CALL,base.relocations.R_BEDROCK_TLSDESC:)
\bottomrule
\end{longtable}
\endgroup
\subsection{Relocation Families and Additional Constraints}
Relocations 1--12 write absolute data, immediate, or effective-address
payloads. Relocations 13--23 write PC-relative, next-instruction-relative, or
section-relative results. Relocations 24--33 address GOT slots, the GOT base,
symbols relative to that base, and PLT entries. Relocations 34--38 are ordinary
dynamic-loader operations. Relocations 39--44 implement GS0 local-exec and
TLSDESC TLS. The table defines widths and calculations; the dynamic-linking and
TLS sections supply additional rules.

\clearpage
\section{Dynamic Linking}\label{page:dynamic-linking}
Dynamic objects use a fixed eager-binding Bedrock GOT and PLT contract.
\subsection{Dynamic Linking Rules}
Ordinary dynamic linking supports standard and indirect functions, copy
relocations, and \texttt{R\_BEDROCK\_IRELATIVE}. The loader resolves every
ordinary \texttt{R\_BEDROCK\_JUMP\_SLOT} and \texttt{R\_BEDROCK\_IRELATIVE}
relocation before transferring control to the program entry point or any
constructor. Shared-object calls use near \texttt{CALL}/\texttt{RET} within
the ABI-visible linkage domain. The GOT base symbol is
\texttt{\_GLOBAL\_OFFSET\_TABLE\_}.

Every ordinary PLT entry occupies 32 bytes and has 16-byte alignment.
\subsection{GOT Reserved Entries}
\BedrockTableCaption{Global Offset Table Reserved Entries}
\begingroup\footnotesize
\setlength{\tabcolsep}{2pt}
\begin{longtable}{@{}p{0.55in}p{1.55in}p{3.4in}@{}}
\toprule
\textbf{Index} & \textbf{Name} & \textbf{Meaning}\\
\midrule
\endhead
0 & \texttt{\_\allowbreak{}DYNAMIC} & dynamic section address\\
1 & reserved & zero\\
2 & reserved & zero\\
\bottomrule
\end{longtable}
\endgroup
\subsection{PLT Layout}
\BedrockTableCaption{Procedure Linkage Table Layout}
\begingroup\footnotesize
\setlength{\tabcolsep}{2pt}
\begin{longtable}{@{}p{0.75in}p{1.55in}p{3.20in}@{}}
\toprule
\textbf{Offset} & \textbf{Bytes or field} & \textbf{Effect}\\
\midrule
\endhead
\texttt{+0x00} & \texttt{e7 c9 80 67} & \texttt{JMP.Q [PC + disp64]}\\
\texttt{+0x04} & \texttt{disp64} & little-endian \texttt{R\_BEDROCK\_GOTPCREL64} relocation field\\
\texttt{+0x0c} & twenty \texttt{01} bytes & canonical \texttt{NOP} padding through offset \texttt{+0x1f}\\
\bottomrule
\end{longtable}
\endgroup

A PLT relocation uses explicit addend \texttt{+4}. Architectural PC names byte zero of the current instruction,
whereas relocation place \texttt{P} is the field at entry offset four; therefore
\texttt{GOT(S)+A-P} evaluates to \texttt{GOT(S)-PLTentry}, the displacement consumed by the PC-relative EA. The
referenced \texttt{.got.plt} slot is 8-byte aligned and contains the eagerly resolved 64-bit near target. Both the
PLT entry and its slot must be reachable under the active CS bounds used by that PC-relative EA.

A PLT transfer preserves \texttt{R0} through \texttt{R7}, \texttt{F0} through
\texttt{F7}, \texttt{GS1} through \texttt{GS5}, \texttt{FLAGS}, and all
segment state. The loader writes each final \texttt{.got.plt} value before
publication. Published slots are immutable and may be covered by
\texttt{PT\_GNU\_RELRO}.
\clearpage
\section{Thread-Local Storage}\label{page:thread-local-storage}
Thread-local storage is addressed through the GS0 segment register. TLS references that remain inside the program and have a link-time constant offset use the local-exec model. TLS references that cross dynamic object boundaries, may be interposed, or otherwise require runtime resolution use the TLSDESC model.
\subsection{TLS Base Model}
GS0 is the TLS base register. A TLS offset is translated by GS0 segment
pre-translation before page-table translation. The local-exec model encodes a
link-time constant byte offset from the GS0 base; TLSDESC supplies runtime
resolution when the reference crosses a dynamic-object boundary or remains
interposable.

\input{abi/elf/documents/fragments/tls_relocation_families.tex}

Taking the address of a TLS object returns an ordinary pre-segment pointer to
that thread instance. Every ABI-permitted GS0 and ordinary-data path consuming
that coordinate resolves it to the same byte of that TLS object. Copying a TLS pointer to
another thread retains its originating thread instance; the pointer becomes invalid when its
originating thread terminates or the defining shared object is unloaded.
\subsection{Model Selection}
\begin{itemize}
\item A TLS symbol defined in the same program image may use local-exec addressing when its byte offset from GS0 is link-time constant.
\item A TLS reference from a shared object or to an interposable TLS symbol uses TLSDESC.
\item The linker may relax a TLSDESC sequence to local-exec when symbol binding and final placement make the GS0-relative offset constant.
\item A source expression \texttt{S+k} uses the descriptor for \texttt{S}; after resolution, the caller adds \texttt{k} to the returned offset. The descriptor relocations encode the base symbol \texttt{S}.
\end{itemize}
\Needspace{1.8in}
\subsection{TLSDESC Entry}
\BedrockTableCaption{TLSDESC Entry Layout}
\begingroup\footnotesize
\setlength{\tabcolsep}{2pt}
\begin{longtable}{@{}p{0.8in}p{1.7in}p{0.65in}p{2.35in}@{}}
\toprule
\textbf{Offset} & \textbf{Field} & \textbf{Type} & \textbf{Meaning}\\
\midrule
\endhead
\texttt{+0x00} & \texttt{function} & \texttt{u64} & near-call address of the resolver or fast function\\
\texttt{+0x08} & \texttt{argument} & \texttt{u64} & signed direct GS0-relative offset or resolver-specific opaque token\\
\bottomrule
\end{longtable}
\endgroup

Each entry occupies 16 bytes in \texttt{.got} and has 16-byte alignment. The
loader resolves the symbol binding, writes both words, and completes any memory
ordering required to publish the pair before application code can observe it.
Both words are immutable after publication and may subsequently be protected by
\texttt{PT\_GNU\_RELRO}. The resolver and fast function may use loader-private
near-call addresses.

The argument word belongs to the implementation of the selected function. A
fast function may load \texttt{[R0+8]} and return it directly. An opaque token
is loader-owned metadata and remains valid while the object that defines
the TLS symbol is loaded. A resolver may perform dynamic per-thread lookup or
allocation. The published descriptor remains immutable for both resolver and
fast-function paths.
\subsection{TLSDESC Call ABI}
The canonical sequence is exactly:
\begin{verbatim}
LEA.Q [PC + descriptor@TLSDESC_GOTPCREL], R0
CALL  [R0]
\end{verbatim}
The \texttt{LEA.Q} places the descriptor address in \texttt{R0}.
\texttt{CALL [R0]} reads the 64-bit near-call target from descriptor offset
zero and preserves \texttt{R0}, so the called function receives the
descriptor address in \texttt{R0}. The canonical sequence places the
\texttt{LEA.Q} and \texttt{CALL} adjacently.

The function returns a valid signed byte offset from the GS0 TLS base in
\texttt{R0}. It preserves \texttt{GS0}, \texttt{R8}--\texttt{R15},
\texttt{F8}--\texttt{F15}, \texttt{V16}--\texttt{V31}, and
\texttt{P8}--\texttt{P15}. The linkage protocol catalog identifies the
resolver scratch set.
Zero is a valid TLS offset. Resolution failure terminates the resolver path.

\subsection{TLSDESC Relocations}
\paragraph{Descriptor address.}
\texttt{R\_BEDROCK\_TLSDESC\_GOTPCREL32S} and
\texttt{R\_BEDROCK\_TLSDESC\_GOTPCREL64} apply to the PC-relative displacement
field of the canonical \texttt{LEA.Q} and compute the address of the descriptor
for \texttt{STT\_TLS} symbol \texttt{S}. The source expression has semantic
TLS addend zero. Under the PC-relative field-bias rule in Section~6.1, the RELA
addend is exactly $A=\delta$ and contains solely the instruction-field offset.

\Needspace{1.0in}
\paragraph{Call marker.}
\texttt{R\_BEDROCK\_TLSDESC\_CALL} has size zero and names the
same TLS symbol as the paired TLSDESC GOTPCREL relocation on the immediately
preceding \texttt{LEA.Q}. Its addend is zero and its place is byte zero of the
exact immediately following \texttt{CALL [R0]}. It is a static-link relaxation
marker consumed exclusively by the static linker. A producer places labels and
control-flow targets outside both instructions of a relaxable sequence.

\paragraph{Descriptor initialization.}
\texttt{R\_BEDROCK\_TLSDESC} applies only to \texttt{STT\_TLS} symbols, has
addend zero, and is placed at offset zero of an aligned descriptor. It directs
the loader to initialize two adjacent 64-bit words with the selected
\texttt{\{function, argument\}} pair for \texttt{S}; this operation initializes
the pair directly. A link or load rejects an unresolved weak TLS symbol because
GS0-relative offset zero represents a valid resolved offset. A weak symbol that resolves to a definition is
handled normally.

The linker may merge descriptors only when they name the same resolved binding
\texttt{S} and have the required zero semantic addend. If any unrelaxed use
remains, a final dynamic link emits one aligned pair in \texttt{.got} and the
corresponding dynamic \texttt{R\_BEDROCK\_TLSDESC}. The loader initializes the
pair eagerly before publication. A fully static executable relaxes every
TLSDESC use; any remaining use makes the static link fail. If all uses are
relaxed, the final link discards the descriptor and its initialization
relocation.

\subsection{TLSDESC Local-Exec Relaxation}
Relaxation is permitted only after final binding proves that \texttt{S} is
non-preemptible and \texttt{TLS(S)} is a link-time constant. It replaces the
entire canonical pair while preserving the section size and every following
instruction address:
\begin{itemize}\raggedright
\item With \texttt{R\_BEDROCK\_TLSDESC\_GOTPCREL32S}, the original
      7-byte \texttt{LEA.Q} and 4-byte \texttt{CALL [R0]} occupy 11 bytes. If
      \texttt{TLS(S)} fits signed 32 bits, the linker emits
      \texttt{LEN 11, MOV.Q TLS(S), R0}, applies
      \texttt{R\_BEDROCK\_TLS\_OFFSET32S} with addend zero to its 32-bit
      immediate payload, and writes four zero padding bytes after the natural
      7-byte encoding.
\item With \texttt{R\_BEDROCK\_TLSDESC\_GOTPCREL64}, the original
      11-byte \texttt{LEA.Q} and 4-byte \texttt{CALL [R0]} occupy 15 bytes. The
      linker emits \texttt{LEN 15, MOV.Q TLS(S), R0}, applies
      \texttt{R\_BEDROCK\_TLS\_OFFSET64} with addend zero to its 64-bit
      immediate payload, and writes four zero padding bytes after the natural
      11-byte encoding.
\end{itemize}
The linker consumes both code relocations. The removed call may eliminate
normal caller-saved clobbers and unobservable resolver stack activity. Code
relies only on the returned offset and defined preserved state. A descriptor
is allocated exactly when an unrelaxed reference remains.

\clearpage
\section{Code Models}\label{page:code-models}
Code models define placement and locality assumptions that assemblers, linkers, and compilers may use when selecting address materialization, displacement, GOT, PLT, and relocation forms. Bedrock 32-bit displacement and direct control-transfer forms are signed, so direct 32-bit relative references cover a -2 GiB to +2 GiB - 1 window.
\BedrockTableCaption{Bedrock Code Models}
\begingroup\footnotesize
\setlength{\tabcolsep}{2pt}
\begin{longtable}{@{}p{0.65in}p{2.25in}p{2.60in}@{}}
\toprule
\textbf{Model} & \textbf{Placement contract} & \textbf{Default access strategy}\\
\midrule
\endhead
\texttt{low} & static addresses in \texttt{0x00000000..0x7fffffff} & signed 32-bit absolute data and direct transfers\\
\texttt{small} & referenced code, data, GOT, and PLT fit one signed 32-bit relative window & PC-relative access; GOT/PLT for interposable symbols\\
\texttt{medium} & code, GOT, and PLT fit one signed 32-bit window; data placement is unrestricted & direct code; GOT-resolved 64-bit data addresses\\
\texttt{high} & static high image with internal references in a signed 32-bit displacement window & image-base or PC-relative access with full-width fallback\\
\texttt{large} & unrestricted 64-bit address space & full-width materialization and indirect transfer\\
\bottomrule
\end{longtable}
\endgroup
\subsection{low Model Details}
The \texttt{low} model is a static image model whose named code and data lie
between \texttt{0x00000000} and \texttt{0x7fffffff}. Code uses direct signed
32-bit branch or call payloads; data uses absolute or base-relative signed
32-bit payloads. The model permits omission of both GOT and PLT.

Its default relocation set is \texttt{R\_BEDROCK\_ABS32S},
\texttt{R\_BEDROCK\_IMM32S}, \texttt{R\_BEDROCK\_DISP32S},
\texttt{R\_BEDROCK\_BRDISP32S}, and \texttt{R\_BEDROCK\_CALL32S}. The linker
may shrink a direct transfer or address materialization after final placement
proves that the result fits a shorter encoding.

\subsection{small Model Details}
The \texttt{small} model is position-independent and assumes that local code,
nearby data, the GOT, and the PLT fit the signed 32-bit relative window of the
use site or selected base. Code uses direct signed 32-bit transfers; data uses
PC-relative or GOTPCREL signed 32-bit materialization.

Its default relocations are \texttt{R\_BEDROCK\_BRDISP32S},
\texttt{R\_BEDROCK\_CALL32S}, \texttt{R\_BEDROCK\_GOTPCREL32S},
\texttt{R\_BEDROCK\_PLT32S}, and
\texttt{R\_BEDROCK\_TLSDESC\_GOTPCREL32S}. The linker may shrink a relative
form or replace GOT/TLSDESC access with a direct form when the resolved symbol
is local and in range.
\subsection{medium Model Details}
The \texttt{medium} model keeps code, GOT, and PLT within a signed 32-bit
relative window but permits data objects anywhere in the 64-bit address space.
Code uses direct signed 32-bit transfers. Data normally uses a GOT-resolved
64-bit address reached through a signed 32-bit GOT reference.

Its default relocations are \texttt{R\_BEDROCK\_BRDISP32S},
\texttt{R\_BEDROCK\_CALL32S}, \texttt{R\_BEDROCK\_GOTPCREL32S},
\texttt{R\_BEDROCK\_GOT64}, \texttt{R\_BEDROCK\_PLT32S}, and
\texttt{R\_BEDROCK\_TLSDESC\_GOTPCREL32S}. A linker may replace a GOT data
reference with a direct relative reference when final placement puts the
object in range.

\subsection{high Model Details}
The \texttt{high} model is a static high-address image. Internal code may use
direct signed 32-bit transfers or image-base-relative addressing. Data uses a
PC-relative GOT or an image-base-relative form when in range; the GOT and PLT
are image-local or base-relative.

Its default relocations are \texttt{R\_BEDROCK\_BRDISP32S},
\texttt{R\_BEDROCK\_CALL32S}, \texttt{R\_BEDROCK\_GOTPCREL32S},
\texttt{R\_BEDROCK\_ABS64}, and \texttt{R\_BEDROCK\_DISP64}. Compact internal
references are permitted when final placement proves the signed 32-bit range;
full-width and base-relative sequences remain valid.

\subsection{large Model Details}
The \texttt{large} model permits arbitrary placement. Code uses full-width
address materialization followed by an indirect transfer when necessary; data,
GOT, and PLT addresses are full width.

Its default relocations are \texttt{R\_BEDROCK\_ABS64},
\texttt{R\_BEDROCK\_PCREL64}, \texttt{R\_BEDROCK\_GOTPCREL64},
\texttt{R\_BEDROCK\_PLT64}, and
\texttt{R\_BEDROCK\_TLSDESC\_GOTPCREL64}. Final placement may justify
relaxing any full-width form to a shorter encoding.

\clearpage
\section{Debugging and Unwind Encoding}\label{page:debugging-and-unwind}
The ELF ABI owns the architecture-wide DWARF register-number assignment used
by every source language and unwind producer.

(:elf-debug-registers:)

An unwind description that saves a vector or predicate register describes its
complete MAX\_VLEN-dependent image. Predicate register locations contain the
packed maximum-width image with its exact architectural byte count.

\clearpage
\section{Assembler Contract}\label{page:assembler-contract}
\BedrockField{Default ABI:}{bedrock-elf}
\subsection{Strict Mode}
\begin{itemize}
\item Pure ISA syntax is accepted.
\item Ambiguous symbolic operands require an explicit relocation annotation.
\item Noncanonical encodings are rejected.
\end{itemize}
\subsection{Relaxation}
\begin{itemize}
\item Branch and call payload widths may be selected automatically when the target is known.
\item Address materialization may shrink from 64-bit to 32-bit, 16-bit, or 8-bit payloads when the relocation result fits.
\item Overlong encodings are emitted only when explicitly requested.
\end{itemize}

\end{document}
